\documentclass[12pt,a4paper,twoside]{book}%      autres choix : book, report
\usepackage[utf8]{inputenc}%           gestion des accents (source)
\usepackage[T1]{fontenc}%              gestion des accents (PDF)
\usepackage[french]{babel}%          gestion du fran�ais
\usepackage{textcomp}%                 caract�res additionnels
\usepackage{mathtools,amssymb}
\usepackage{amsthm}
\usepackage{xurl}
\usepackage{mathrsfs}
\usepackage{amsmath}% packages de l'AMS + mathtools
\usepackage{lmodern}
\usepackage[bottom=2.5cm,left=2.5cm,right=2.5cm,top=2.5cm]{geometry}
\usepackage[colorlinks=true, linkcolor=blue, citecolor=blue, urlcolor=blue]{hyperref}
\usepackage{fancyhdr}

\usepackage{graphicx}
\usepackage{xcolor}
\pagestyle{fancy}
\fancyhf{}
\theoremstyle{definition}
\newtheorem{dfn}{Définition}[section]
\newtheorem*{pr}{\textbf{Preuve}}
\newtheorem{rem}{\textit{Remarque}}[section]
\newtheorem{thm}[dfn]{\textit{Théorème}}
\newtheorem{prop}[dfn]{\textit{Proposition}}
\newtheorem{cor}[dfn]{\textit{Corollaire}}
\newtheorem{lem}[dfn]{\textit{Lemme}}
\newtheorem{exm}[dfn]{\textit{Exemple}}
\newtheorem{prp}[dfn]{\textbf{Propriété}}
\newtheorem{for}[dfn]{\textit{Formule}}
\numberwithin{equation}{section}
\newtheorem{ex}{Exemple}[subsection]
%\setcounter{page}{1}
\title{ \textbf{Théorème Central Limite et Principe de Déviations Modérées pour les Équations Différentielles Stochastiques de McKean-Vlasov)}}
\begin{document}
	\maketitle
	\newpage
	\tableofcontents
	\newpage
	
	
		\chapter{Rappel et Préliminaire}
Soient \( |\cdot| \) et \( \langle \cdot, \cdot \rangle \) respectivement la norme euclidienne et le produit scalaire dans \( \mathbb{R}^d \)
	\section{Quelque définition des espaces et distance}
	\begin{dfn}
	Soit $T>0$ fixé. On note $\Omega = C([0,T],\mathbb{R}^d)$ l'espace des fonction continues $h : [0,T] \to\mathbb{R}^d$ telle que $h(0) = 0.$
	
	L'espace de Cameron-Martin, noté $\mathbb{H}$, est l'ensemble des fonction $h\in \Omega$ qui sont absolument continues et dont la dérivée (au sens de Lebesgue) est de carrée intégrable:
	\[
	\mathbb{H} = \big\{h\in C([0,T];\mathbb{R}^d):h(0)=\textbf{0},\dot{h}(t) \text{existe p.p }t ,\|h\|_H :=\Big(\int_{0}^{T}|\dot{h}(t)|^2dt\Big)^{\frac{1}{2}} \big\}
	\] où \textbf{0} désigne le vecteur dont toutes les composantes sont nulles.
	
	On munit $\mathbb{H}$ du produit scalaire
	\[
	\langle h,k \rangle_{\mathbb{H}} := \int_{0}^{T}\dot{h}(t)\cdot \dot{k}(t)dt 
	\]
	Muni de cette structure, $\mathbb{H}$ est un espace de Hilbert séparable.
	\end{dfn}
	\begin{dfn}
		Un espace métrique $X$ est une espace Polonais de $X$ est séparable et il existe une distance complète induisant la topologie de $X$.
	\end{dfn}
	\begin{dfn}
		Soit $\mathcal{P}(\mathbb{R}^d)$ l'espace des mesures de probabilité sur $\mathbb{R}^d$.\\
		On note $\mathcal{P}_2(\mathbb{R}^d)$ est le sous-ensemble des mesure de probabilité sur $\mathbb{R}^d$ possédant un moment d'ordre deux fini, c'est-à-dire:
		\[
		\mathcal{P}_2(\mathbb{R}^d) = \left\{ \mu \text{ probabilité sur }\mathbb{R}^d : \mu(|\cdot|^2) = \int_{\mathbb{R}^d} |x|^2\mu(dx)<\infty
		\right\} 
		\]
		où $\mu(f):=\int fd\mu$ pour un fonction mesurable $f$
	\end{dfn}
	\begin{rem}
		Le moment d'ordre 2 garantit que l'on peut définir des distances entre mesures via la distance de Wasserstein $W_2$
	\end{rem}
	\begin{dfn}
		Soient $\mu$ et $\nu\in\mathcal{P}_2(\mathbb{R}^d)$, la distance de Wasserstein d'ordre 2 entre $\mu$ et $\nu$ est définie par la formule:
		\[
		W_2(\mu, \nu) := \left( \inf_{\pi \in \mathcal{C}(\mu, \nu)} \int_{\mathbb{R}^d \times \mathbb{R}^d} |x - y|^2 \, \pi(dx, dy) \right)^{1/2}, 
		\quad \mu, \nu \in \mathcal{P}_2(\mathbb{R}^d),
		\]
		où $\mathcal{C}(\mu,\nu)$ désigne l'ensemble des couplages de $\mu$ et $\nu$; c'est-à-dire que $\pi\in\mathcal{C}(\mu,\nu)$ est une mesure de probabilité sur $\mathbb{R}^d\times\mathbb{R}^d$ telle que $\pi(\cdot\times\mathbb{R}^d) = \mu$ et $\pi(\mathbb{R}^d\times\cdot)=\nu$.
	\end{dfn}
	\begin{prop}
		L'espace de mesure de probabilité $\mathcal{P}_2(\mathbb{R}^d)$ est une espace polonais pour distance Wasserstein
	\end{prop}
	\begin{pr}
		à recherche
	\end{pr}
	
	%\begin{dfn}
%	Pour toute $\mu\in\mathcal{P}_2(\mathbb{R}^d)$, l'espace tangente en $\mu$ est donné par:

%	\[T_{\mu,2}=L^{2}(\mathbb{R}^{d}\to\mathbb{R}^{d};\mu):=\{\phi:\mathbb{R}^{d}\to\mathbb{R}^{d}\text{ est mesurable avec }\mu(|\phi|^{2})<\infty\}		\]
%	\end{dfn}
%Pour $\phi\in T_{\mu,2}$
%\[\|\phi\|_{T_{\mu,2}}^{2}=\int_{\mathbb{R}^{d}}|\phi(x)|^{2}\mu(\mathrm{d}x).\] est la norme canonique de l'espace $T_{\mu,2}$
\section{Équation différentielle stochastique de type McKean-Vlasov}
\subsection{Équation différentielle stochastiques classique}
\begin{dfn}
	Une équation différentielle stochastique (SDE) sur $\mathbb{R}^d$ est une équation de la forme:
	\begin{equation}
		dX_t = b(t,X_t)dt + \sigma(t,X_t)dB_t, X_0 = x_0\in\mathbb{R}^d
	\end{equation}
	où:
	\begin{itemize}
		\item $B_t$ est un mouvement brownien standard de dimension $d$.
		\item $b:[0,T]\times\mathbb{R}^d\longrightarrow\mathbb{R}^d$ est la dérivée drift
		\item $\sigma:[0,T]\times\mathbb{R}^{d\times d}$ est la diffusion.
	\end{itemize}
\end{dfn}
\begin{thm}[Théorème d’existence et d’unicité pour les équations différentielles stochastiques]
	Soit $T > 0$ et soient 
	\[
	b : [0,T] \times \mathbb{R}^n \to \mathbb{R}^n, 
	\quad 
	\sigma : [0,T] \times \mathbb{R}^n \to \mathbb{R}^{n \times m}
	\]
	des fonctions mesurables satisfaisant
	\begin{equation} \label{eq:lin_growth}
		|b(t,x)| + |\sigma(t,x)| \leq C(1+|x|), 
		\quad \forall (t,x) \in [0,T]\times \mathbb{R}^n,
	\end{equation}
	pour une certaine constante $C$, où 
	\[
	|\sigma(t,x)|^2 = \sum_{i,j} \sigma_{ij}^2(t,x).
	\]
	
	De plus, supposons que
	\begin{equation} \label{eq:lipschitz}
		|b(t,x)-b(t,y)| + |\sigma(t,x)-\sigma(t,y)| \leq D |x-y|, 
		\quad \forall x,y \in \mathbb{R}^n, \, t \in [0,T],
	\end{equation}
	pour une certaine constante $D$.
	
	Soit $Z$ une variable aléatoire indépendante de la $\sigma$-algèbre 
	$\mathcal{F}^{(m)}$ engendrée par $(B_s)_{s \geq 0}$, et telle que
	\[
	\mathbb{E}[x_0^2] < \infty.
	\]
	
	Alors, l’équation différentielle stochastique
	\[
	dX_t = b(t,X_t)\,dt + \sigma(t,X_t)\,dB_t, 
	\quad 0 \leq t \leq T, 
	\quad X_0 = x_0,
	\]
	admet une solution unique $X_t(\cdot)$, continue en $t$, et vérifiant :
	\begin{equation}
		X_t(\cdot) \text{ est adaptée à la filtration } 
		\mathcal{F}^{x_0}_t \text{ engendrée par } x_0 \text{ et } (B_s)_{s \leq t},
	\end{equation}
	et
	\begin{equation} 
		\mathbb{E}\!\left[\int_0^T |X_t|^2 \, dt\right] < \infty.
	\end{equation}
\end{thm}

\subsection{Équation différentielle stochastique de type McKean-Vlasov (MV-SDE)}
Dans cette section, nous rappelons l’équation différentielle stochastique de McKean–Vlasov (MV-SDE), définie sur l’espace euclidien $(\mathbb{R}^d, \langle \cdot, \cdot \rangle ,|\cdot|)$. 
Soit $(\Omega,\mathcal{F}_t,\{\mathcal{F}_t\}_{t\geqslant 0},\mathbb{P})$ un espace de probabilité filtré complet et $\mathcal{L}_{X^\varepsilon_t}$ désigne la loi du processus stochastique $X^\varepsilon_t$.
\begin{dfn}
	Une équation différentielle stochastique de type Mckean-Vlasov (MV-SDE) sur $(\mathbb{R}^d, \langle \cdot, \cdot \rangle ,|\cdot|)$ est de la forme:
	\begin{equation}
		dX_t^\varepsilon = b_t(X_t^\varepsilon,\mathcal{L}(X_t^\varepsilon))dt +\sqrt{\varepsilon}\sigma_t(X_t^\varepsilon,\mathcal{L}(X_t^\varepsilon))dW_t,\text{    } X_0^\varepsilon=x
	\end{equation}
	où $\varepsilon > 0$ est appelé le para
	mètre d'échelle.Ici, $W_t$ est une mouvement brownien standard $d-$dimensionnel.
\end{dfn}
\subsubsection{Hypothèses sur les coefficients}
\textbf{(H1)} Les coefficients $b : [0, \infty) \times \mathbb{R}^d \times \mathcal{P}_2(\mathbb{R}^d) \to \mathbb{R}^d$ et $\sigma : [0, \infty) \times \mathbb{R}^d \times \mathcal{P}_2(\mathbb{R}^d) \to \mathbb{R}^{d \otimes d}$ sont continus. Il existe une fonction croissante $K : [0, \infty) \to [0, \infty)$ telle que
\begin{equation}
	\max\left\{ \|\nabla b_t(\cdot, \mu)(x)\|, \| D^L b_t(x, \cdot)(\mu) \|_{T_{\mu,2}} \right\} \leq K(t), \quad t \geq 0,\ x \in \mathbb{R}^d,\ \mu \in \mathcal{P}_2(\mathbb{R}^d),
\end{equation}
\begin{equation}
	\|\sigma_t(x, \mu) - \sigma_t(y, \nu)\| \leq K(t)(|x - y| + \mathbb{W}_2(\mu, \nu)), \quad t \geq 0,\ x, y \in \mathbb{R}^d,\ \mu, \nu \in \mathcal{P}_2(\mathbb{R}^d),
\end{equation}
et
\begin{equation}
	|b_t(0, \delta_0)| + \|\sigma_t(0, \delta_0)\| \leq K(t), \quad t \geq 0.
\end{equation}

\textbf{(H2)} Le coefficient $b_t(x, \mu)$ est différentiable par rapport à $x$ et $\mu$, respectivement, et les fonctions dérivées satisfont
\begin{equation}
	\|\nabla b_t(., \mu)(x) - \nabla b_t(., \nu)(y)\| \leq K(t)(|x - y| + \mathbb{W}_2(\mu, \nu)),
\end{equation}
\[
\begin{aligned}
	&\left| \mathbb{E}\left[ \langle D^L b_t(x, \cdot)(\mathcal{L}_X)(X), \phi \rangle \right] 
	- \mathbb{E}\left[ \langle D^L b_t(y, \cdot)(\mathcal{L}_Y)(Y), \phi \rangle \right] \right| \\
	&\leq K(t)\left( |x - y| + \mathbb{W}_2(\mathcal{L}_X, \mathcal{L}_Y) + \left( \mathbb{E}[|X - Y|^2] \right)^{1/2} \right) \left( \mathbb{E}[|\phi|^2] \right)^{1/2}
\end{aligned}
\]

pour tout $t \geq 0$, $x, y \in \mathbb{R}^d$, $\mu, \nu \in \mathcal{P}_2(\mathbb{R}^d)$ et $X, Y, \phi \in L^2(\Omega \to \mathbb{R}^d, \mathbb{P}$
\begin{rem}
	D'après (H1), on a pour tout $t\geq0$ $x,y\in\mathbb{R}^d,\mu,\nu\in\mathcal{P}_2(\mathbb{R}^d)$ que:
	\begin{equation}
		|b_t(x, \mu) - b_t(y, \nu)| \leq K(t)(|x - y| + \mathbb{W}_2(\mu, \nu)).
	\end{equation}
\end{rem}

Intuitivement, lorsque le paramétré $\varepsilon$ tend vers 0 dans (1.2.6), le terme de diffusion disparaître et on obtient l'équation différentielle ordinaire suivante:
\begin{equation}
	dX_t^0=b_t(X^0_t,\delta_{X_t^0})dt.
\end{equation}
avec la même donnée initiale que (1.2.6), c'est-à-dire $X_0^0=x$. Comme $x$ est déterministe, on déduit que $\delta_{X^0_\cdot}$ est une mesure de Dirac centrée sur la trajectoire $X^0$
\section{Dérivée de Lions}
Pour tout $\mu\in\mathcal{P}_2(\mathbb{R}^d)$, on définit l'espace tangent en $\mu$ par
\[
T_{\mu,2} = L^2(\mathbb{R})^d\longrightarrow \mathbb{R}^d :=\left\{\phi:\mathbb{R}^d\longrightarrow\mathbb{R}^d \text{ est mesurable avec }\int_{\mathbb{R}^d}|\phi(x)|^2\mu(dx)<\infty\right\}
\]
la norme associé est donnée par:
\[
\|\phi\|^2_{T_{\mu,2}}:=\int_{\mathbb{R}^d}|\phi(x)|^2\mu(dx)\] 

\begin{dfn}
	Soit $f:\mathcal{P}_2(\mathbb{R}^d)\longrightarrow\mathbb{R}^d$ un fonction continue, et $I_d$ la carte identité sur $\mathbb{R}^d$
	
	\begin{enumerate}
		\item Si, pour tout $\mu\in\mathcal{P}_2(\mathbb{R}^d)$, on a
		\[
		T_{\mu,2}\ni\phi \mapsto D^L_{\phi}f(\mu):=\lim_{\varepsilon\to 0}\frac{f(\mu\circ(I_d+\varepsilon\phi)^{-1})-f(\mu)}{\varepsilon}\in\mathbb{R}
		\]
		et que cette limite définit un fonctionnel linéaire borné, alors on dit que $f$ est intrinsèquement différentiable en $\mu$, et on note $D^L_{\phi}f(\mu)$ la dérivée intrinsèque de $f$ en $\mu$

	\item Si, pour tout $\mu\in\mathcal{R}_2(\mathbb{R})$, on a
	\[
\lim_{\|\phi\|_{T_{\mu,2}} \to 0} \frac{|f(\mu \circ (I_d + \phi)^{-1}) - f(\mu) - D^L_\phi f(\mu)|}{\|\phi\|_{T_{\mu,2}}} = 0.
 \]
 alors on dit que $f$ est $L-$différentiable en $\mu$, et l'on note $D_\mu f(\mu)$ la $L-$dérivée (ou dérivée de Lions) de $f$ en $\mu$
\\ Dans ce cas, par le théorème de représentation de Riez, l'élément unique $D^Lf(\mu)\in T_{\mu,2}$ satisfaisant
 \[
 \langle D^L f(\mu), \phi \rangle_{T_{\mu,2}} := \int_{\mathbb{R}^d} \langle D^L f(\mu)(x), \phi(x) \rangle \mu(dx) = D^L_\phi f(\mu), \quad \forall \phi \in T_{\mu,2},
 \]
 \item  On écrit \( f \in C^1(\mathcal{P}_2(\mathbb{R}^d)) \) si \( f \) est \( L \)-différentiable en tout point \( \mu \in \mathcal{P}_2(\mathbb{R}^d) \), et la dérivée \( L \)-dérivée admet version \( D^L f(\mu)(x) \) jointe continu en \( (x, \mu) \in \mathbb{R}^d \times \mathcal{P}_2(\mathbb{R}^d) \). Si de plus, \( D^L f(\mu)(x) \) est bornée, on dénote \( f \in C_b^1(\mathcal{P}_2(\mathbb{R}^d)) \).
 
 Pour une fonction à valeurs vectorielles $f = (f_i)$, ou une fonction à valeurs matricielles $f = (f_{ij})$ avec des composantes $L$-différentiables, on écrit
 \[
 D^L_\phi f(\mu) = (D^L_\phi f_i(\mu)), \quad \text{ou} \quad D^L_\phi f(\mu) = (D^L_\phi f_{ij}(\mu)), \quad \mu \in \mathcal{P}_2(\mathbb{R}^d).
 \] 
 	\end{enumerate}
\end{dfn}


\section{Principe de grandes déviation (PGD) et Principe de déviation modéré (PDM)}
\subsection{Principe de grandes déviation (PGD)}
Le principe des grandes déviation (PGD) décrit le comportement asymptotique d'une famille de mesures de probabilité $\{\mu_\varepsilon\}_{\varepsilon>0}$ définit sur un espace topologique $(X,\mathcal{B}_X)$, lorsque $\varepsilon \to 0$, au moyen d'une fonction de taux, à travers des bornes exponentielles supérieures et inférieures portants sur les probabilités des ensembles mesurables de $X$.
\begin{dfn}
	Soit $X$ un espace polonais muni de la tribu borélienne $\mathcal{B}$.  
	On dit qu’une fonction $I : X \to [0, \infty]$ est une \emph{fonction de taux} (rate function) si :  
	\begin{enumerate}
		\item $I$ n’est pas identiquement égale à $+\infty$ ;
		\item $I$ est semi-continue inférieurement ;
		\item pour tout $L \geq 0$, l’ensemble $\{x \in X : I(x) \leq L\}$ est compact.
	\end{enumerate}
\end{dfn}

\begin{dfn}
	Une famille $\{\mu_\varepsilon\}$ de mesures de probabilité sur $(X, \mathcal{B})$ 
	est dite satisfaire le \textbf{principe des grandes déviations} (PGD) 
	de fonction de taux $I$ si:
	\begin{equation}
		\varlimsup_{\varepsilon \downarrow 0} \varepsilon\log\mu_\varepsilon (F)\leq -\inf_{x\in F}I(x)
	\end{equation}
	pour tout ensemble fermé $F$ de $X$ et
	\begin{equation}
	\varliminf_{\varepsilon \downarrow 0} \varepsilon\log\mu_\varepsilon(G)\geq -\inf_{x\in G}
	\end{equation}
	pour tout ensemble ouvert non vide $G$ de $X$  
	\end{dfn}

\section{Théorème centrale limite (TCL)}







\end{document}